\chapter{Drinking Alcohol}

\section*{Definition and Overview}
Alcohol is a legal depressant drug that affects the brain and body, and drinking it is deeply embedded in Australian social life. While moderate drinking is common, alcohol is a leading contributor to illness, injury, and death in Australia. The National Health and Medical Research Council (NHMRC) guidelines recommend that healthy adults drink no more than ten standard drinks a week and no more than four on any one day, and that the safest option for some groups, including pregnant women and people under 18, is not to drink at all. This chapter explains the effects of alcohol, the risks, and how to drink responsibly.

\section*{Epidemiology}
Alcohol is the most widely used drug in Australia, with around three-quarters of adults drinking. Around one in four Australians drinks at levels that put their health at risk, and one in four children under 14 is exposed to risky drinking by a parent. Alcohol contributes to around 5,000 deaths and more hospitalisations in Australia each year than tobacco, and it is a factor in violence, road crashes, and suicide. Alcohol use is higher among men and varies with age, with young adults the heaviest drinkers.

\section*{Aetiology and Pathophysiology}
\begin{itemize}
    \item \textbf{Effects on the brain:} alcohol slows brain activity, impairing judgement, coordination, and reaction time
    \item \textbf{Standard drinks:} a standard drink contains 10 grams of alcohol, equivalent to a middy of beer, a small glass of wine, or a nip of spirits
    \item \textbf{Metabolism:} the liver breaks down alcohol at a fixed rate, and excess accumulates in the blood
    \item \textbf{Blood alcohol concentration:} depends on the amount drunk, body size, and time
    \item \textbf{Short-term risks:} injury, accidents, violence, and risky behaviour
    \item \textbf{Long-term risks:} liver disease, cancer, heart disease, brain damage, and dependence
\end{itemize}

\section*{Clinical Features}
\begin{itemize}
    \item The NHMRC guidelines: no more than ten standard drinks per week and four per day for healthy adults
    \item Drinking less is better for health; there is no level that is completely safe
    \item Alcohol-free days help reduce the weekly total
    \item Pregnant women, people under 18, and people with certain conditions should not drink
    \item Signs of intoxication: slurred speech, poor coordination, and impaired judgement
    \item Harm from alcohol can occur even below the guidelines
\end{itemize}

\section*{Differential Diagnosis}
\begin{itemize}
    \item Low-risk, risky, and high-risk drinking, which need different responses
    \item Alcohol dependence versus harmful use
    \item Binge drinking versus daily drinking, which have different risks
    \item Alcohol-related illness versus other causes of the same symptoms
\end{itemize}

\section*{When to Seek Medical Care}
People who are concerned about their drinking, who drink above the guidelines, or who have health or relationship problems related to alcohol should discuss it with their GP. Withdrawal from heavy drinking should be medically supervised.

\textbf{Call 000 (or local emergency services) for:}
\begin{itemize}
    \item Signs of alcohol poisoning: confusion, vomiting, seizures, slow breathing, or unconsciousness
    \item A person who cannot be roused after drinking
    \item Withdrawal symptoms such as seizures, hallucinations, or severe confusion
    \item Injury or violence related to drinking
\end{itemize}

\section*{Diagnosis and Investigations}
\begin{itemize}
    \item \textbf{Assessment of drinking:} review of how much, how often, and the effects
    \item \textbf{Screening tools:} questionnaires such as the AUDIT identify risky drinking
    \item \textbf{Health assessment:} blood pressure, liver function, and other tests
    \item \textbf{Assessment of dependence:} withdrawal symptoms, tolerance, and loss of control
    \item \textbf{Mental health review:} depression, anxiety, and other factors
\end{itemize}

\section*{Management}
\begin{itemize}
    \item \textbf{Follow the guidelines:} no more than ten drinks weekly and four daily
    \item \textbf{Plan alcohol-free days:} reduces the weekly total
    \item \textbf{Pace drinking:} sip drinks, eat before and during drinking, and alternate with water
    \item \textbf{Use standard measures:} count standard drinks rather than glasses
    \item \textbf{Avoid binge drinking:} no more than four standard drinks on one occasion
    \item \textbf{Seek support:} counselling, apps, and programs for cutting down
    \item \textbf{Treatment for dependence:} medical withdrawal support and ongoing treatment
\end{itemize}

\section*{Complications}
\begin{itemize}
    \item Injury, road crashes, falls, and violence
    \item Alcohol poisoning and overdose
    \item Liver disease, including cirrhosis
    \item Cancers of the mouth, throat, liver, and breast
    \item Heart disease, high blood pressure, and stroke
    \item Dependence and addiction
    \item Harm to unborn babies from drinking in pregnancy
\end{itemize}

\section*{Prognosis and Outcomes}
Reducing alcohol intake improves health quickly, and people who drink within the guidelines have lower risks of injury, disease, and premature death. Even heavy drinkers who cut down or stop see improvements in blood pressure, liver health, mood, and sleep. Dependence is treatable, and many people recover with support. Drinking less is one of the most effective steps anyone can take for better health.

\section*{Prevention and Self-Care}
\begin{itemize}
    \item Know the guidelines and count your standard drinks
    \item Have alcohol-free days each week
    \item Eat before and while drinking
    \item Alternate alcoholic drinks with water
    \item Never drink and drive
    \item Do not drink if pregnant, trying to conceive, or breastfeeding
    \item Keep alcohol out of reach of children and model responsible drinking
    \item Seek help early if drinking is hard to control
\end{itemize}

\section*{Sources and Further Reading}
The following reputable sources provide further information for healthcare students and patients seeking reliable, current advice about drinking alcohol.

\begin{itemize}
    \item Alcohol and Drug Foundation. \textit{Alcohol: information and support.} https://adf.org.au/
    \item National Health and Medical Research Council. \textit{Australian guidelines to reduce health risks from drinking alcohol.} https://www.nhmrc.gov.au/
    \item DrinkWise Australia. \textit{Responsible drinking information.} https://drinkwise.org.au/
    \item Healthdirect Australia. \textit{Alcohol: effects and guidelines.} https://www.healthdirect.gov.au/
    \item NHS. \textit{Alcohol units and guidelines.} https://www.nhs.uk/live-well/
    \item World Health Organization. \textit{Alcohol fact sheet.} https://www.who.int/
\end{itemize}
